\subsection{受控格式的一致性、收敛与结构}
\label{sec:controlled_convergence}

加入控制后，首先区分力幅值与状态增量。对本标量磁场端口，步后修正为$U_k=(0,h u_kd\cos(d^TX))$，其整体量级随$h$缩小；力幅值$u_k$本身不必随步长趋零。以$m=2$的标量有势系统作精确参考检验，省略$h$的固定增量负对照在步长从$0.2$减至$0.0125$时，终点误差由$0.3365$增至$6.7919$，而正确加载的I/G与III呈约二阶、II趋近一阶。该负对照在固定外生输入下仍可保辛，说明保辛不能替代一致性与收敛。

沿用同一两模态梁和固定控制律，将原步长依次取$4h_0,2h_0,h_0,h_0/2,h_0/4$，分别考察三种参考。第一种为光滑外生输入的连续流，各方法在时间中点评价输入；第二种为实际步起点反馈保持与各自连续反馈流的差，包含采样和积分误差；第三种先从连续反馈参考在步起点采样并冻结整条输入序列，让三种格式共享同一序列，再与该保持序列的连续分段流比较，以分离积分误差。第三种是外生序列诊断，并非另一个实际反馈闭环。

\begin{table}[t]
\centering\small\setlength{\tabcolsep}{4pt}
\caption{加入控制后的观察阶。梁模型采用$2\times2$质量矩阵；外生输入行为零初值，其余两行为10个已训练模型与一个无补偿基线。梁按全时域相对$L^2$、标量按终点$L^2$计算最细两档观察阶；保留四位小数后各模型取值一致。}
\label{tab:controlled_convergence}
\begin{tabular}{lrrr}
\toprule
参考定义 & I/G & II & III \\
\midrule
光滑外生输入 & 2.0000 & 1.0121 & 2.0000 \\
冻结保持序列 & 2.0000 & 1.0121 & 2.0000 \\
实际采样反馈 & 1.0068 & 1.0028 & 1.0068 \\
标量$m=2$、外生 & 1.9999 & 0.9927 & 1.9999 \\
\bottomrule
\end{tabular}
\end{table}


零初值结果见表\ref{tab:controlled_convergence}，预先固定的三个非零测试初值结果也全部保留。冻结序列的保持流参考经二倍子步加密，最大状态变化小于$9.4\times10^{-15}$。I/G与III的步内积分约二阶，却不消除实际采样反馈相对连续反馈的近一阶总误差。G已由原生成函数标量控制项独立核对与I等价，因而不作为独立统计重复。

在各步长下进一步区分固定$u_k$与完整反馈Jacobian：前者辛缺陷最大约$5.7\times10^{-16}$，后者计入$D_x\Phi+\Phi_uD\kappa$后最大约$0.0211$，与独立状态差分的最大差小于$2.8\times10^{-11}$。反馈缺陷随步长减小不意味着每个固定步长映射严格保辛；上述观察阶也仅对应已定义的系统、参考和输入保持方式。
